\en{In this chapter, we develop the terms and propositions about the Weierstraß elliptic functions that we need for our proof of the Chudnovsky formula. The notation and some proofs are based on \cite{\en{EN}Freitag2000}, where one can find more detailed explanations.}%
\de{In diesem Kapitel werden alle für den Beweis der Chudnovsky-Formel benötigten Begriffe und Sätze entwickelt. Die Notation und einige Beweise orientieren sich an \cite{\en{EN}Freitag2000}. Dort findet man auch ausführlichere Erläuterungen und Zusammenhänge.}

\begin{defi}\label{\en{EN}defgitter} \en{For each pair $(\omega_1,\omega_2)$ of complex numbers which is $\mathbb R$-linearly independent (which means $\omega_2/\omega_1\notin\mathbb R$) we call}%
\de{Zu jedem Paar $(\omega_1,\omega_2)$ komplexer Zahlen, die $\mathbb R$-linear unabhängig sind (also gilt $\omega_2/\omega_1\notin\mathbb R$) nennt man}
$$L = \mathbb Z \omega_1 + \mathbb Z \omega_2 = \left\{m\omega_1+n\omega_2~|~m,n\in\mathbb Z\right\}\subset\mathbb C$$
\en{a "lattice". $\omega_1$ and $\omega_2$ are called "basic periods" of the lattice.}%
\de{ein \glqq Gitter\grqq. $\omega_1$ und $\omega_2$ heißen dann auch \glqq Basisperioden\grqq~des Gitters.}
\end{defi}

\begin{defi}
\en{An "elliptic function" is a meromorphic function $f:\mathbb C \rightarrow \mathbb C \cup \{\infty\}$ with the property}%
\de{Eine \glqq elliptische Funktion zum Gitter $L$\grqq~ist eine meromorphe Funktion $f:\mathbb C \rightarrow \mathbb C \cup \{\infty\}$ mit der Eigenschaft} $$f(z+\omega)=f(z)\qquad\text{\en{for all}\de{für alle} }~\omega\in L~\text{ \en{and}\de{und} }~z\in\mathbb C$$
\en{"Meromorphic" means that $f$ has no essential singularities, that the set of poles of $f$ has no accumulation point, and that $f$ is holomorphic apart from the poles.}%
\de{\glqq Meromorph\grqq~bedeutet, dass $f$ keine außerwesentlichen Singularitäten hat, dass die Polstellenmenge von $f$ keinen Häufungspunkt hat, und dass $f$ außerhalb der Polstellen analytisch ist.}
\en{To show that a meromorphic function is elliptic it suffices to check if $f(z+\omega_1)=f(z)=f(z+\omega_2)$ holds for both basic periods of the lattice -- that's why elliptic functions are also called "doubly periodic".}%
\de{Um nachzuweisen, dass eine meromorphe Funktion elliptisch ist, reicht es zu prüfen, ob $f(z+\omega_1)=f(z)=f(z+\omega_2)$ für die beiden Basisperioden des Gitters gilt -- deshalb nennt man elliptische Funktionen auch \glqq doppeltperiodisch\grqq.}
\end{defi}

\begin{defi}\label{\en{EN}defimodz}
\en{Every lattice $L$ produces an equivalence relation on the complex numbers: We call $z_1\in \mathbb C$ and $z_2\in \mathbb C$ "equivalent modulo $L$", iff it holds $z_1-z_2\in L$ (since it doesn't matter if one uses $z_1$ or $z_2$ in an elliptic function of the lattice $L$).}%
\de{Jedes Gitter $L$ erzeugt eine Äquivalenzrelation auf den komplexen Zahlen: Wir nennen $z_1\in \mathbb C$ und $z_2\in \mathbb C$ \glqq äquivalent modulo $L$\grqq, falls $z_1-z_2\in L$ ist (dann ist es nämlich egal, ob man $z_1$ oder $z_2$ in eine elliptische Funktion zum Gitter $L$ einsetzt).}
\end{defi}

\begin{defi}\label{\en{EN}fund}
\en{The "fundamental parallelogram" $\mathcal P$ and its closure $\overline{\mathcal P}$ are:}%
\de{Das \glqq Periodenparallelogramm\grqq~$\mathcal P$ und sein Abschluss $\overline{\mathcal P}$ lauten:}
$$\mathcal P = \left\{~s\omega_1+t\omega_2~|~0\leq s,t< 1~\right\}
\qquad\text{\en{and}\de{und}}\qquad
\overline{\mathcal P}= \left\{~s\omega_1+t\omega_2~|~0\leq s,t\leq 1~\right\}$$
\en{Since $\omega_1$ and $\omega_2$ are $\mathbb R$-linearly independent, it holds: for all $z\in\mathbb C$ there is exactly one $z'\in\mathcal P$ which is equivalent to $z$ (modulo $L$).
Figure~\ref{\en{EN}abbwege} on p.~\pageref{\en{EN}abbwege} depicts $\overline{\mathcal P}$.}%
\de{Da $\omega_1$ und $\omega_2$ $\mathbb R$-linear unabhängig sind, gibt es zu jedem $z\in\mathbb C$ genau ein $z'\in\mathcal P$, das modulo $L$ äquivalent zu $z$ ist.
Eine Abbildung von $\overline{\mathcal P}$ befindet sich auf Seite~\pageref{\en{EN}abbwege} (Abb.~\ref{\en{EN}abbwege}).}
\end{defi}

\begin{thm}[\en{Liouville's Theorem}\de{Satz von Liouville}]\label{\en{EN}liouville0}
\en{Any bounded analytic function $\mathbb C \rightarrow \mathbb C$ is constant.}%
\de{Jede beschränkte analytische Funktion $\mathbb C \rightarrow \mathbb C$ ist konstant.}
\end{thm}
\begin{proof}
\en{Given $z\in \mathbb C$, we will prove $f'(z)=0$: By deriving Cauchy's integral formula (with Leibniz's rule) we obtain for all $r>0$:}%
\de{Für $z\in \mathbb C$ beweisen wir zunächst $f'(z)=0$: Durch Ableiten der Cauchy'schen Integralformel (mit der Leibniz'schen Regel) folgt für alle $r>0$:}
\begin{align*}
    |f'(z)| &= \left|\frac{1}{2\pi i} \oint_{|\zeta-z|=r}\frac{f(\zeta)}{(\zeta-z)^2}d\zeta\right|
\leq \frac{1}{2\pi} \cdot \frac{C}{r^2}\cdot 2\pi r = \frac{C}{r}
\end{align*}
\en{Here we used the boundedness $|f(\zeta)|\leq C$ and the perimeter of the circle.}%
\de{Hier haben wir die Beschränktheit $|f(\zeta)|\leq C$ und den Umfang des Kreises genutzt.}
\en{For $r\rightarrow\infty$ we obtain $f'(z)=0$ -- thus $f$ is constant.}%
\de{Für $r\rightarrow\infty$ folgt also $f'(z)=0$ und somit dass $f$ konstant ist.}
\end{proof}

\begin{thm}[\en{First Liouville Theorem}\de{Erster Liouville'scher Satz}]\label{\en{EN}liouville1}
\en{Any elliptic function without poles is constant.}%
\de{Jede elliptische Funktion ohne Polstellen ist konstant.}
\end{thm}
\begin{proof}
\en{Any elliptic function $f$ with basic periods $\omega_1$ and $\omega_2$ takes any of its values in the fundamental parallelogram $\mathcal P$ (cf.~Def.~\ref{\en{EN}fund}). But its closure $\overline{\mathcal P}$ is closed and bounded (see Fig.~\ref{\en{EN}abbwege} on p.~\pageref{\en{EN}abbwege}). Since $f$ has no poles, $|f|$ is continuous and must have a maximum in $\overline{\mathcal P}$. But then, because of its periodicity, $f$ is bounded on the whole complex plane. From Liouville's theorem (Prop.~\ref{\en{EN}liouville0}) we deduce that $f$ must be constant.}%
\de{Eine elliptische Funktion $f$ nimmt jeden ihrer Werte schon im Periodenparallelogramm $\mathcal P$ an. Aber $\overline{\mathcal P}$ ist beschränkt und abgeschlossen (siehe Abb.~\ref{\en{EN}abbwege} auf S.~\pageref{\en{EN}abbwege}), deshalb besitzt $f$ wie jede stetige Funktion in $\overline{\mathcal P}$ ein Maximum. Aus der Periodizität folgt, dass $f$ auf ganz $\mathbb C$ beschränkt und somit nach dem Satz von Liouville (Satz~\ref{\en{EN}liouville0}) konstant ist.}
\end{proof}

\begin{thm}[\en{Second Liouville Theorem}\de{Zweiter Liouville'scher Satz}]\label{\en{EN}liouville2}
\en{Any elliptic function has only finitely many poles (modulo $L$) and the sum of their residues vanishes.}%
\de{Jede elliptische Funktion hat nur endlich viele Pole (modulo $L$) und die Summe ihrer Rediduen verschwindet.}
\end{thm}
\begin{proof}
\en{For any pole of an elliptic function, there is an equivalent pole in $\mathcal P$ (cf.~Def.~\ref{\en{EN}fund}). The set of poles of an elliptic function is discrete, thus only finitely many poles are in the closure $\overline{\mathcal P}$ of the fundamental parallelogram ($\overline{\mathcal P}$ is compact).}%
\de{Zu jeder Polstelle einer elliptischen Funktion gibt es eine äquivalente Polstelle in $\mathcal P$ (vgl.~Def.~\ref{\en{EN}fund}). Die Menge der Pole einer elliptischen Funktion ist diskret, also liegen nur endlich viele Pole in  der kompakten Menge $\overline{\mathcal P}$ aus Def.~\ref{\en{EN}fund}.
}
\en{Then we move $\overline{\mathcal P}$ so that no more poles are on its border, and integrate along the border. Since $\wp$ is doubly periodic, this integral vanishes (since the integrals along parallel parts of the border cancel each other out). From the residue theorem we deduce that the sum of the residues vanishes.}%
\de{Wenn wir nun $\overline{\mathcal P}$ so verschieben, dass keine Pole mehr auf dem Rand liegen, und einmal entlang des Randes integrieren, folgt aus der Periodizität der $\wp$-Funktion, dass das Integral den Wert Null hat (die Beiträge gegenüberliegender Kanten heben sich gegenseitig auf). Aus dem Residuensatz folgt nun, dass die Summe der Residuen verschwindet.}
\end{proof}

\begin{thm}[\en{Third Liouville Theorem}\de{Dritter Liouville'scher Satz}]\label{\en{EN}liouville3}
\en{Any non-constant elliptic function $f$ has the same number of zeros and poles modulo $L$, if they are counted with their multiplicities.}%
\de{Jede nichtkonstante elliptische Funktion $f$ hat modulo $L$ gleich viele Null- und Polstellen, wobei diese mit ihrer Vielfachheit zu rechnen sind.}
\end{thm}
\begin{proof}
\en{If $f$ is a non-constant elliptic function, $g(z):=\frac{f'(z)}{f(z)}$ also is a non-constant elliptic function.}%
\de{Wenn $f$ eine nichtkonstante elliptische Funktion ist, ist auch $g(z):=\frac{f'(z)}{f(z)}$ eine nichtkonstante elliptische Funktion.}

\en{If the Laurent series of $f$ in $z_0$ starts with $f(z)\approx a\cdot (z-z_0)^{k}$ (where $k\in \mathbb Z$, $k\neq 0$), it holds $f'(z)\approx k\cdot a\cdot(z-z_0)^{k-1}$ and $g(z)\approx \frac{k}{z-z_0}$. Thus every pole and every zero of $f(z)$ produces a pole of $g(z)$ of order one with residue $k$.}%
\de{Wenn die Laurentreihe von $f$ in $z_0$ mit $f(z)\approx a\cdot (z-z_0)^{k}$ beginnt ($k\in \mathbb Z$, $k\neq 0$), gilt $f'(z)\approx k\cdot a\cdot(z-z_0)^{k-1}$ und $g(z)\approx \frac{k}{z-z_0}$. Also gilt für alle Null- und Polstellen von $f(z)$, dass $g(z)$ dort einen Pol erster Ordnung mit Residuum $k$ hat.}

\en{From its definition $g(z):=\frac{f'(z)}{f(z)}$ we see that $g$ has no further poles.}%
\de{Aus der Definition $g(z):=\frac{f'(z)}{f(z)}$ folgt, dass $g$ keine weiteren Pole hat.}

\en{The sum of the residues of $g$ vanishes (second Liouville theorem, Prop.~\ref{\en{EN}liouville2}), thus we have:
The sum of the positive residues of $g$ (the sum of the multiplicities of the zeros of $f$) has the same absolute value as the sum of the negative residues of $g$ (the sum of the multiplicities of the poles of $f$).}%
\de{Die Summe der Residuen von $g$ verschwindet (zweiter Liouville'scher Satz~\ref{\en{EN}liouville2}), also gilt:
Die Summe der positiven Residuen von $g$ (die Summe der Nullstellenordnungen von $f$) ist betragsmäßig gleich groß wie die Summe der negativen Residuen von $g$ (die Summe der Polstellenordnungen von $f$).}
\end{proof}


\begin{defi}\label{\en{EN}defisigma}
\en{The Weierstraß $\sigma$-function of the lattice $L$ is defined as follows:}%
\de{Die Weierstraß'sche $\sigma$-Funktion zum Gitter $L$ ist wie folgt definiert:}
$$\sigma(z;L):=z\cdot\prod_{\substack{\omega\in L\\\omega\neq 0}}\left\{\left(1-\frac z \omega \right)\cdot\exp\lk\frac z \omega + \frac 1 2 \left(\frac z \omega\right)^2\rk\right\}$$
\en{The $\sigma$-function will be analyzed further in chapter~\ref{\en{EN}kapFourier}, see for example Prop.~\ref{\en{EN}sigmaodd},~\ref{\en{EN}trafosigma} and~\ref{\en{EN}fouriersigma}.}%
\de{Die $\sigma$-Funktion wird in Kapitel~\ref{\en{EN}kapFourier} weiter untersucht, siehe z.B.~Satz~\ref{\en{EN}sigmaodd},~\ref{\en{EN}trafosigma} und~\ref{\en{EN}fouriersigma}.}
\end{defi}

\begin{bem}\label{\en{EN}bemsigma}
\en{This product converges absolutely because of the exponential factor, and the zeros of $\sigma(z;L)$ are exactly the points of the lattice $L$ and are zeros of order $1$. Nevertheless, the $\sigma$-function is \emph{not} doubly periodic (cf.~Prop.~\ref{\en{EN}trafosigma}).}%
\de{Dieses Produkt konvergiert aufgrund des Exponentialfaktors absolut, und die Nullstellen von $\sigma(z;L)$ sind genau die Punkte des Gitters $L$. Es sind einfache Nullstellen. Trotzdem ist die $\sigma$-Funktion \emph{nicht} doppeltperiodisch (vgl.~Satz~\ref{\en{EN}trafosigma}).}
\end{bem}

\begin{defi}\label{\en{EN}defizeta}
\en{The Weierstraß $\zeta$-function of a lattice $L$ is defined as the logarithmic derivative of the $\sigma$-function, whose product yields a sum because of~~$\ln(a\cdot b)=\ln a +\ln b$:}%
\de{Die Weierstraß'sche $\zeta$-Funktion zu einem Gitter $L$ ist als logarithmische Ableitung der $\sigma$-Funktion definiert:}%, wobei das Produkt wegen $\ln(a\cdot b)=\ln a +\ln b$ in eine Summe übergeht:}
\begin{align*}
\zeta(z;L) &:= \frac {d}{dz}\ln\sigma(z;L) = \frac {d}{dz}\lk\ln z\rk + \sum_{\substack{\omega\in L\\\omega\neq 0}}\frac {d}{dz}\left\{\ln\left(1-\frac z \omega \right)+\frac z \omega + \frac 1 2 \left(\frac z \omega\right)^2\right\}\\
&= \frac 1 z + \sum_{\substack{\omega\in L\\ \omega\neq 0}} \left(\frac 1 {z-\omega} +\frac 1 \omega +\frac z {\omega^2}\right)
\end{align*}
\en{The $\zeta$-function will be analyzed further in chapter~\ref{\en{EN}kapquasiint}, see for example Def.~\ref{\en{EN}defetak} and Rem.~\ref{\en{EN}bempitch}.}%
\de{Die $\zeta$-Funktion wird in Kapitel~\ref{\en{EN}kapquasiint} weiter untersucht, siehe z.B.~Def.~\ref{\en{EN}defetak} und Bem.~\ref{\en{EN}bempitch}.}
\end{defi}

\begin{defi}\label{\en{EN}defiwp}
\en{The Weierstraß $\wp$-function denotes the negative derivative of the Weierstraß $\zeta$-function:}%
\de{Die Weierstraß'sche $\wp$-Funktion ist definiert als die negative Ableitung der Weierstraß'schen $\zeta$-Funktion:}
$$\wp(z;L) :=  -\zeta'(z;L) = \frac 1 {z^2} + \sum_{\substack{\omega\in L\\ \omega\neq 0}} \left(\frac 1 {(z-\omega)^2} -\frac 1 {\omega^2}\right)$$
\end{defi}

\begin{bem}\label{\en{EN}pstrich}
\en{The derivative of the Weierstraß $\wp$-function reads:}%
\de{Die Ableitung der Weierstraß'schen $\wp$-Funktion lautet:}
$$\wp'(z;L)=\sum_{\omega\in L}\frac{-2}{(z-\omega)^3}$$
\end{bem}

\begin{thm}\label{\en{EN}pgerade}
\en{$\wp(z;L)$ is an even function and $\wp'(z;L)$ is an odd function, i.e.}%
\de{$\wp(z;L)$ ist eine gerade und $\wp'(z;L)$ ist eine ungerade Funktion, d.h.}
$$\wp(-z;L)=\wp(z;L)\qquad\text{\en{and}\de{und}}\qquad\wp'(-z;L)=-\wp'(z;L)$$
\end{thm}
\begin{proof}
\en{If $\omega$ runs through all points of the lattice, then $-\omega$ does it too:}%
\de{Mit $\omega$ durchläuft auch $-\omega$ alle Gitterpunkte. Hieraus folgt:}
\begin{align*}
    \wp(-z;L) &= \frac 1 {(-z)^2} + \sum_{\substack{\omega\in L\\ \omega\neq 0}} \left(\frac 1 {(-z-\omega)^2} -\frac 1 {\omega^2}\right)\\
            &= \frac 1 {z^2} + \sum_{\substack{-\omega\in L\\ -\omega\neq 0}} \left(\frac 1 {(z-(-\omega))^2} -\frac 1 {(-\omega)^2}\right) = \wp(z;L)
\end{align*}
\en{And for $\wp'(z)$ it holds:}%
\de{Und für $\wp'(z)$ gilt:}\belowdisplayskip=-12pt
\begin{align*}
    \wp'(-z;L) &= \sum_{\omega\in L}\frac{-2}{(-z-\omega)^3} = - \sum_{-\omega\in L}\frac{-2}{(z-(-\omega))^3} = -\wp'(z;L)
\end{align*}
\end{proof}

\begin{thm}\label{\en{EN}wpdoppeltper}
\en{The Weierstraß $\wp$-function is doubly periodic, i.e.~for all $\omega\in L$ we have $\wp(z+\omega;L)=\wp(z;L)$.}%
\de{Die Weierstraß'sche $\wp$-Funktion ist doppeltperiodisch, d.h.~für alle $\omega\in L$ gilt $\wp(z+\omega;L)=\wp(z;L)$.}
\end{thm}
\begin{proof}
\en{$\wp'$ is doubly periodic, since the summation runs through all lattice points and since there are no further terms (see Remark~\ref{\en{EN}pstrich}).
So we get $\wp'(z+\omega)-\wp'(z)=0$ and thus $\wp(z+\omega)-\wp(z)=\text{const}$.
If $\omega$ is a basic period of the lattice, then $-\frac \omega 2 \notin L$.
We get the value of the constant with Prop.~\ref{\en{EN}pgerade}: $\wp\lk-\frac \omega 2 + \omega\rk - \wp\lk-\frac \omega 2\rk =\wp\lk\frac \omega 2\rk - \wp\lk-\frac \omega 2\rk = 0$.
This yields $\wp(z+\omega)=\wp(z)$ for all \emph{basic} periods of the lattice $L$ and thus for \emph{all} points of the lattice.}%
\de{Man sieht sofort, dass $\wp'$ doppeltperiodisch ist, weil über alle Gitterpunkte summiert wird und keine weiteren Terme in der Summe stehen.
Also gilt $\wp'(z+\omega)-\wp'(z)=0$ und somit $\wp(z+\omega)-\wp(z)=\text{const}$.
Wenn wir für $\omega$ eine der Basisperioden des Gitters wählen, dann ist $-\frac \omega 2 \notin L$.
Wir erhalten den Wert der Konstanten mit Hilfe von Satz~\ref{\en{EN}pgerade}: $\wp\lk-\frac \omega 2 + \omega\rk - \wp\lk-\frac \omega 2\rk =\wp\lk\frac \omega 2\rk - \wp\lk-\frac \omega 2\rk = 0$.
Es folgt $\wp(z+\omega)=\wp(z)$ für alle Basisperioden des Gitters $L$ und somit auch für alle anderen Gitterpunkte.}
\end{proof}

\begin{thm}\label{\en{EN}zerowp}
\en{The zeros of $\wp'$ are exactly those points $\frac\omega 2$, for which $\omega\in L$ but $\frac\omega 2\notin L$.
If $L=\mathbb Z \omega_1 + \mathbb Z \omega_2$, this yields the following three zeros (see Fig.~\ref{\en{EN}abbwege} on p.~\pageref{\en{EN}abbwege}):}%
\de{Die Nullstellen von $\wp'$ sind genau diejenigen Stellen $\frac\omega 2$, die selbst nicht im Gitter liegen, für die aber $\omega$ im Gitter liegt.
Wenn $L=\mathbb Z \omega_1 + \mathbb Z \omega_2$ ist, gilt also}
$$\wp'\lk\frac{\omega_1}{2}\rk = \wp'\lk\frac{\omega_2}{2}\rk = \wp'\lk\frac{\omega_1+\omega_2}{2}\rk=0$$
\end{thm}
\begin{proof}
\en{Choose $\omega_k\in L$ so that $\frac{\omega_k}2\notin L$. Then we get: if $\omega$ runs through all points of the lattice, then also $\omega'=\omega+\omega_k$ does it. This yields:}%
\de{Sei $\omega_k\in L$ so gewählt, dass $\frac{\omega_k}2\notin L$. Dann durchläuft mit $\omega$ auch $\omega'=\omega+\omega_k$ alle Gitterpunkte und es gilt:}
\begin{align*}
    \wp'\lk-\frac{\omega_k}2;L\rk &= \sum_{\omega\in L}\frac{-2}{\left(-\frac{\omega_k}2-\omega\right)^3} = \sum_{\omega'\in L}\frac{-2}{\left(-\frac{\omega_k}2-(\omega'-\omega_k)\right)^3}\\
    &= \sum_{\omega'\in L}\frac{-2}{\left(\frac{\omega_k}2-\omega'\right)^3} = \wp'\lk\frac{\omega_k}2;L\rk
\end{align*}
\en{From Prop.~\ref{\en{EN}pgerade} we know that $\wp'$ is \emph{odd} and according to our premises $\pm\frac{\omega_k}2$ is not in $L$. From this we get $\wp'\lk-\frac{\omega_k}2;L\rk = -\wp'\lk\frac{\omega_k}2;L\rk$.
This yields $\wp'\lk\frac{\omega_k}2;L\rk=0$.}%
\de{Weil $\wp'$ nach Satz~\ref{\en{EN}pgerade} eine \emph{ungerade} Funktion ist und weil nach Voraussetzung $\pm\frac{\omega_k}2\notin L$ ist, gilt $\wp'\lk-\frac{\omega_k}2;L\rk = -\wp'\lk\frac{\omega_k}2;L\rk$.
Folglich muss $\wp'\lk\frac{\omega_k}2;L\rk=0$ sein.}
\en{Using the third Liouville theorem (Prop.~\ref{\en{EN}liouville3}) we see that $\wp'$ has no further zeros (modulo $L$).}%
\de{Aus dem dritten Liouville'schen Satz~\ref{\en{EN}liouville3} folgt, dass $\wp'$ (modulo $L$) keine weiteren Nullstellen hat.}
\end{proof}

\begin{defi}\label{\en{EN}defeisen}
\en{The series}\de{Die Reihen} $\displaystyle G_n = G_n(L) := \sum_{\substack{\omega\in L\\ \omega\neq 0}} \omega^{-n}$
\en{are called "Eisenstein series of weight $n$" and converge absolutely for natural $n\geq 3$.}%
\de{heißen \glqq Eisensteinreihen zum Gitter $L$\grqq\ und konvergieren für natürliche $n\geq 3$ absolut.}
\end{defi}

\begin{thm}\label{\en{EN}eisenungerade}
\en{The Eisenstein series of odd weight vanish (i.e.~they take on the value $0$).}%
\de{Die Eisensteinreihen mit ungeradem Gewicht verschwinden.}
\end{thm}
\begin{proof}
\en{Since $n$ is odd, we can deduce that for all $\omega\in L-\{0\}$ the summands $\omega^{-n}$ and $(-\omega)^{-n}=-\left(\omega^{-n}\right)$ cancel each other out. Thus the full sum takes on the value $0$.}%
\de{Wenn $n$ ungerade ist, dann gilt für alle $\omega\in L-\{0\}$, dass sich die Summanden $\omega^{-n}$ und $(-\omega)^{-n}=-\left(\omega^{-n}\right)$ gegenseitig aufheben.}
\end{proof}

\begin{thm}\label{\en{EN}laurentwp}
\en{The Weierstraß $\wp$-function admits the following Laurent series expansion around $z=0$ without a constant term:}%
\de{Die Weierstraß'sche $\wp$-Funktion lässt sich in der Nähe von $z=0$ durch folgende Laurentreihe ohne konstanten Term darstellen:}
\begin{align*}
\wp(z;L)&=\frac 1 {z^2} + \sum_{n=1}^\infty (2n+1)\cdot G_{2n+2}(L)\cdot z^{2n}
\end{align*}
\end{thm}
\begin{proof}
\en{First we analyze $f(z):=\wp(z;L)-\frac 1{z^2}$. From Def.~\ref{\en{EN}defiwp} we get $f(0)=0$.
Then we get the derivatives of $f(z)$ at $z=0$ with the representation of $\wp'$ from Remark~\ref{\en{EN}pstrich}:}%
\de{Wir untersuchen zunächst $f(z):=\wp(z;L)-\frac 1{z^2}$. Bei $z=0$ folgt direkt aus der Definition~\ref{\en{EN}defiwp} der $\wp$-Funktion, dass $f(0)=0$ ist.
Für die Ableitungen von $f(z)$ bei $z=0$ folgt dann mit Hilfe der Darstellung von $\wp'$ aus Bemerkung~\ref{\en{EN}pstrich}:}
$$f^{(n)}(z) = (-1)^n(n+1)!\sum_{\substack{\omega\in L\\ \omega\neq 0}} \frac{1}{(z-\omega)^{n+2}}\qquad\text{\en{if}\de{falls} }n\geq 1$$
\pagebreak\\
\en{From Prop.~\ref{\en{EN}eisenungerade} we deduce that the odd derivatives vanish at $z=0$, and that the even derivatives are:}%
\de{Hieraus folgt (wegen Satz~\ref{\en{EN}eisenungerade}), dass die ungeraden Ableitungen bei $z=0$ verschwinden, und dass für die geraden gilt:}
$$f^{(2n)}(0) = (-1)^{2n}(2n+1)!\sum_{\substack{\omega\in L\\ \omega\neq 0}} \frac{1}{(-\omega)^{2n+2}}=(2n+1)!\sum_{\substack{\omega\in L\\ \omega\neq 0}} \frac{1}{\omega^{2n+2}}=(2n+1)!\cdot G_{2n+2}$$
\en{with the Eisenstein series from Definition~\ref{\en{EN}defeisen}.
Thus we have shown that it holds $f(z)=\sum_{n=1}^\infty \frac{f^{2n}(0)}{(2n)!}\cdot z^{2n} = \sum_{n=1}^\infty (2n+1)G_{2n+2}\cdot z^{2n}$ and the proposition is proven.}%
\de{mit den Eisensteinreihen aus Definition~\ref{\en{EN}defeisen}.
Insgesamt haben wir also bewiesen, dass $f(z)=\sum_{n=1}^\infty \frac{f^{2n}(0)}{(2n)!}\cdot z^{2n} = \sum_{n=1}^\infty (2n+1)G_{2n+2}\cdot z^{2n}$ gilt, und wir sind fertig.}
\end{proof}

\begin{thm}\label{\en{EN}dglP}
\en{The Weierstraß $\wp$-function satisfies the algebraic differential equation:}%
\de{Es gilt folgende algebraische Differentialgleichung der $\wp$-Funktion zum Gitter $L$:}
\begin{align*}
    \wp'(z)^2 &= 4 \wp(z)^3 - g_2 \wp(z) - g_3\\
    \text{\en{with}\de{mit}}\quad g_2 &= g_2(L) := 60 G_4(L) = 60 \sum_{\substack{\omega\in L\\ \omega\neq 0}} \omega^{-4}\\
    \text{\en{and}\de{und}}\quad g_3 &= g_3(L) := 140 G_6(L) = 140 \sum_{\substack{\omega\in L\\ \omega\neq 0}} \omega^{-6}
\end{align*}
\end{thm}
\begin{proof}
\en{We use the beginning of the Laurent series expansion from Prop.~\ref{\en{EN}laurentwp} and show that $h(z):=\wp'(z)^2-4\wp(z)^3+60G_4\wp(z)$ has no poles:}%
\de{Wir verwenden den Anfang der Laurentreihe aus Satz~\ref{\en{EN}laurentwp} und zeigen, dass die Funktion $h(z):=\wp'(z)^2-4\wp(z)^3+60G_4\wp(z)$ keine Pole hat:}
\begin{align*}
    \wp(z;L) &= z^{-2} + 3 G_4 z^2 + 5 G_6 z^4 + O(z^6)\\
    \Longrightarrow\quad\wp(z;L)^2 &= z^{-4} + 6 G_4 + 10 G_6 z^2 + O(z^4)\\
    \Longrightarrow\quad\wp(z;L)^3 &= \wp(z;L)^2\cdot\wp(z;L) = z^{-6} + 9 G_4 z^{-2} + 15 G_6 + O(z^2)\\
    \text{\en{and}\de{und}}\quad\wp'(z;L) &= -2z^{-3} + 6 G_4 z + 20 G_6 z^3 + O(z^5)\\
    \Longrightarrow\quad\wp'(z;L)^2 &= 4z^{-6} -24 G_4 z^{-2} - 80 G_6 + O(z^2)\\
    \Longrightarrow\wp'(z;L)^2-4\wp(z;L)^3 &= -60 G_4 z^{-2} - 140 G_6 + O(z^2)\\
    \Longrightarrow\wp'(z;L)^2 -4\wp(z;L)^3 &+ 60 G_4 \wp(z;L) = - 140 G_6 + O(z^2)
\end{align*}
\en{Here we recognize that $h(z)$ has no pole at $z=0$. From the definition of $h(z)$ we see that it is doubly periodic, and that $h(z)$ has no poles in any lattice points. Since neither $\wp$ nor $\wp'$ have poles besides the lattice points we deduce that $h(z)$ is an elliptic function without poles.
Thus, by the first Liouville theorem (Prop.~\ref{\en{EN}liouville1}), $h(z)$ is constant. The value of this constant is $-140G_6$ (see above), so we get
$\wp'(z;L)^2 = 4\wp(z;L)^3 - 60 G_4(L)\wp(z;L) - 140 G_6(L)$.}%
\de{$h(z)$ hat also bei $z=0$ keinen Pol. Weil $h(z)$ aufgrund seiner Definition doppeltperiodisch ist, hat $h(z)$ also auch in den anderen Gitterpunkten keine Pole, und außerhalb der Gitterpunkte haben sowieso weder $\wp$ noch $\wp'$ Pole, insofern ist $h(z)$ eine elliptische Funktion ohne Pole und nach dem ersten Liouville'schen Satz~\ref{\en{EN}liouville1} konstant. Der konstante Wert von $h(z)$ ergibt sich zu $-140G_6$ und wir erhalten $\wp'(z;L)^2 = 4\wp(z;L)^3 - 60 G_4(L)\wp(z;L) - 140 G_6(L)$, was zu zeigen war.}
\end{proof}


\begin{thm}\label{\en{EN}pstrichprod}
\en{If $L=\mathbb Z \omega_1 + \mathbb Z \omega_2$, then it holds:}%
\de{Wenn $L=\mathbb Z \omega_1 + \mathbb Z \omega_2$ ist, gilt:} 
\begin{align*}
    \wp'(z)^2&=4\cdot \left(\wp(z)-e_1\right) \cdot\left(\wp(z)-e_2\right) \cdot\left(\wp(z)-e_3\right)\\
    \intertext{\en{with the pairwise distinct half lattice values of the $\wp$-function}\de{mit den paarweise verschiedenen Halbwerten der $\wp$-Funktion}}
    e_1 &:= \wp\lk\frac{\omega_1}{2}\rk;\qquad e_2 := \wp\lk\frac{\omega_2}{2}\rk;\qquad e_3 := \wp\lk\frac{\omega_1+\omega_2}{2}\rk
\end{align*}
\end{thm}
\begin{proof}
\en{Prop.~\ref{\en{EN}zerowp} tells (for example) $\wp'\lk\frac{\omega_1}{2}\rk=0$.
If we set $f(z):=\wp(z)-e_1$, we get both $f\lk\frac{\omega_1}{2}\rk=0$ and $f'\lk\frac{\omega_1}{2}\rk=0$ -- thus $f$ has a double zero at $\frac{\omega_1}{2}$.
Now the third Liouville theorem (Prop.~\ref{\en{EN}liouville3}) tells that $\wp(z)-e_1$ has no further zeros.
Thus the $e_{1;2;3}$ are pairwise distinct.}%
\de{Nach Satz~\ref{\en{EN}zerowp} gilt z.B.~$\wp'\lk\frac{\omega_1}{2}\rk=0$.
Also gilt für $f(z):=\wp(z)-e_1$ sowohl $f\lk\frac{\omega_1}{2}\rk=0$ als auch $f'\lk\frac{\omega_1}{2}\rk=0$ -- somit hat $f$ bei $\frac{\omega_1}{2}$ eine doppelte Nullstelle.
Aus dem dritten Liouville'schen Satz~\ref{\en{EN}liouville3} folgt, dass $\wp(z)-e_1$ keine weiteren Nullstellen hat.
Somit sind die $e_{1;2;3}$ paarweise verschieden.}

\en{From Prop.~\ref{\en{EN}dglP} and~\ref{\en{EN}zerowp} we deduce that $P(X):=4X^3-g_2X-g_3$ has the three distinct zeros $e_1$, $e_2$ and $e_3$. This proves $P(X)=4(X-e_1)(X-e_2)(X-e_3)$.

Using $X=\wp(z)$ and Prop.~\ref{\en{EN}dglP} proves the Proposition.}%
\de{Aus Satz~\ref{\en{EN}dglP} und~\ref{\en{EN}zerowp} folgt dann, dass $P(X):=4X^3-g_2X-g_3$ die drei verschiedenen Nullstellen $e_1$, $e_2$ und $e_3$ hat. Hieraus folgt $P(X)=4(X-e_1)(X-e_2)(X-e_3)$.
Mit $X=\wp(z)$ folgt dann aus Satz~\ref{\en{EN}dglP} die zu beweisende Aussage.}
\end{proof}
